\chapter{Conical Shock}

A conical shock forms when supersonic flow encounters a cone. Because the body is axisymmetric, the shock wave itself is also a cone anchored at the cone tip. After crossing the shock, the gas undergoes continuous compression as it flows along curved streamlines toward the cone surface. This makes the conical shock problem fundamentally different from the oblique shock: the downstream flow properties cannot be found from simple algebraic relations and instead must be obtained by integrating the Taylor--Maccoll ordinary differential equation from the shock surface down to the cone surface. The cone angle $\theta_c$ and surface Mach number $M_c$ are not inputs to the ODE --- they emerge as outputs of the integration. This uniquely reversed structure means that the shock angle $\theta_s$ is the true central variable of the problem. The mathematical formulations and reference Python scripts for the conical shock numerical integrations were provided by PhD scholar Mr. Mohankumar Neelakandan.

\section{Nomenclature}

\begin{center}
\begin{tabular}{@{}cl@{}}
\toprule
Symbol & Description \\
\midrule
$\gamma$ & Specific heat ratio \\
$M_1$ & Freestream Mach number \\
$\theta_s$ & Shock wave angle \\
$\theta_c$ & Cone half-angle (body angle) \\
$M_c$ & Mach number at the cone surface \\
$\delta$ & Flow deflection angle immediately behind the shock \\
$M_2$ & Mach number immediately behind the shock \\
$\frac{T_2}{T_1}$ & Static temperature ratio immediately behind the shock \\
$\frac{P_2}{P_1}$ & Static pressure ratio immediately behind the shock \\
$\frac{\rho_2}{\rho_1}$ & Density ratio immediately behind the shock \\
$\frac{P_{02}}{P_{01}}$ & Stagnation pressure ratio across the shock \\
$\frac{T_c}{T_1}$ & Static temperature ratio at the cone surface \\
$\frac{P_c}{P_1}$ & Static pressure ratio at the cone surface \\
$\frac{\rho_c}{\rho_1}$ & Density ratio at the cone surface \\
$C_p$ & Surface pressure coefficient \\
$\frac{\Delta s}{R}$ & Non-dimensional entropy rise across the shock \\
$V_r'$ & Radial velocity component, normalized by $V_{max}$ \\
$V_\theta'$ & Tangential velocity component, normalized by $V_{max}$ \\
$V_{total}'$ & Total velocity component, normalized by $V_{max}$ \\
\bottomrule
\end{tabular}
\end{center}

\section{Governing Idea}

In supersonic flow over a cone, the shock wave itself is a cone. When gas passes through this shock, it first turns sharply. After that, it flows along curved paths (streamlines) toward the solid cone surface, compressing gradually. Because this compression is continuous and curved, we cannot calculate the flow properties using simple equations. Instead, we must solve a differential equation called the Taylor--Maccoll equation. Whichever parameter the user inputs, the solver always calculates the shock angle $\theta_s$ first, then integrates the flow equations from that shock angle down to the cone surface where the flow runs parallel to the cone.

\section{Mathematical Flow and Solver Steps}

Analyzing conical flow requires solving the governing differential relations across the shock and along the conical streamlines, and finding the limits of attached flow. The solver executes in the following sequence:

\subsection*{Step 0: Mandatory Setup and Input Mapping}

The specific heat ratio $\gamma > 1$ and the freestream Mach number $M_1 > 1$ are mandatory variables. The user must supply one additional input: the shock angle $\theta_s$, the cone angle $\theta_c$, or the cone surface Mach number $M_c$. Whichever input is given, the solver's goal is to find or verify the shock angle $\theta_s$, which serves as the starting point for all downstream calculations.

\subsection*{Step 1: Finding Attached Shock Limits}

Before starting the main calculation, the solver determines the physical boundaries. Once $M_1$ is inputted, the engine sweeps shock angles $\theta_s$ from the Mach angle $\mu_1$ up to $90^\circ$ to find where the shock detaches from the cone:
\begin{equation}
\mu_1 = \sin^{-1}\left(\frac{1}{M_1}\right)
\end{equation}
By executing the central solver recursively across the range $[\mu_1,\ 90^\circ]$, the engine determines:
\begin{enumerate}[nosep]
    \item The maximum possible cone angle $\theta_{c,max}$ (limit for attached shock).
    \item The minimum possible surface Mach number $M_{c,min}$.
\end{enumerate}

\subsection*{Step 2: Core Solver Dispatching Logic}

Once input limits are set, the solver chooses one of two paths based on the active field:
\begin{itemize}
    \item \textbf{Direct Path:} If the user inputs the shock angle $\theta_s$ (Case 1), the solver validates the range and invokes the central solver directly.
    \item \textbf{Iterative Path:} If the user inputs the cone angle $\theta_c$ (Case 2) or surface Mach number $M_c$ (Case 3), the solver uses Brent's method (root-finding) to solve for the corresponding shock angle $\theta_s$ in the search range $[\mu_1,\ 90^\circ]$. Once $\theta_s$ is found, the engine runs the central solver a final time.
\end{itemize}

\subsection*{Step 3: The Central Solver (Taylor--Maccoll ODE Integration)}

Starting from a known shock angle $\theta_s$, this core solver computes all properties using a forward integration sequence:
\begin{enumerate}
    \item \textbf{Oblique-Shock Jump Conditions:} Applies the exact oblique-shock relations right behind the shock:
    \begin{equation}
    \delta = \tan^{-1}\left[\frac{2}{\tan\theta_s}\times\frac{M_1^2\sin^2\theta_s-1}{M_1^2(\gamma+\cos2\theta_s)+2}\right]
    \end{equation}
    \begin{equation}
    M_2 = \frac{1}{\sin(\theta_s-\delta)}\sqrt{\frac{\frac{\gamma-1}{2}M_1^2\sin^2\theta_s+1}{\gamma M_1^2\sin^2\theta_s-\frac{\gamma-1}{2}}}
    \end{equation}
    \begin{equation}
    \frac{P_2}{P_1} = \frac{2\gamma M_1^2\sin^2\theta_s-(\gamma-1)}{\gamma+1}
    \end{equation}
    \begin{equation}
    \frac{\rho_2}{\rho_1} = \frac{(\gamma+1)M_1^2\sin^2\theta_s}{(\gamma-1)M_1^2\sin^2\theta_s+2}
    \end{equation}
    \begin{equation}
    \frac{T_2}{T_1} = \frac{P_2/P_1}{\rho_2/\rho_1}
    \end{equation}
    If $M_2$ is not a valid real, non-negative value, the trial is rejected.

    \item \textbf{Taylor--Maccoll Initial Conditions:} The velocity behind the shock is scaled as $V' = V/V_{max}$ and decomposed into radial and tangential components ($V_r', V_\theta'$):
    \begin{equation}
    V' = \left[\frac{2}{(\gamma-1)M_2^2}+1\right]^{-1/2}
    \end{equation}
    \begin{equation}
    V_r'(\theta_s) = V'\cos(\theta_s-\delta), \qquad V_\theta'(\theta_s) = -V'\sin(\theta_s-\delta)
    \end{equation}

    \item \textbf{Runge-Kutta-Fehlberg 4(5) (RK45) Integration:} The differential equations are solved step-by-step from the shock to the cone surface. The RK45 method is an adaptive integration scheme. At each step, it calculates both a 4th-order and a 5th-order estimate of the velocity. The difference between these two estimates gives the numerical error. If the error is too large, the step size is automatically reduced; if the error is tiny, the step size is increased. This ensures high accuracy (with a relative tolerance of $10^{-13}$ and absolute tolerance of $10^{-15}$) while keeping the calculation fast. The integrated system is:
    \begin{equation}
    \frac{dV_r'}{d\theta} = V_\theta', \qquad
    \frac{dV_\theta'}{d\theta} = \frac{V_r'V_\theta'^2 - \dfrac{\gamma-1}{2}(1-V_r'^2-V_\theta'^2)\left(2V_r'+V_\theta'\cot\theta\right)}{\dfrac{\gamma-1}{2}(1-V_r'^2-V_\theta'^2) - V_\theta'^2}
    \end{equation}

    \item \textbf{Solid Cone Surface Detection:} The integration stops when the tangential velocity component crosses zero ($V_\theta' = 0$). The crossing angle $\theta_c$ is located with Brent's method root-finder to a tolerance of $10^{-14}$. At this point:
    \begin{equation}
    V_c' = V_r'(\theta_c), \qquad M_c = \sqrt{\frac{2V_c'^2}{(\gamma-1)(1-V_c'^2)}}
    \end{equation}

    \item \textbf{Post-Processing and Ratios Evaluation:} Evaluates the remaining outputs:
    \begin{equation}
    \frac{P_{02}}{P_{01}} = \frac{P_2}{P_1}\left[\frac{1+\frac{\gamma-1}{2}M_2^2}{1+\frac{\gamma-1}{2}M_1^2}\right]^{\frac{\gamma}{\gamma-1}}
    \end{equation}
    \begin{equation}
    \frac{T_c}{T_1} = \frac{T_2}{T_1}\times\frac{1+\frac{\gamma-1}{2}M_2^2}{1+\frac{\gamma-1}{2}M_c^2}
    \end{equation}
    \begin{equation}
    \frac{P_c}{P_1} = \frac{P_2}{P_1}\left[\frac{1+\frac{\gamma-1}{2}M_2^2}{1+\frac{\gamma-1}{2}M_c^2}\right]^{\frac{\gamma}{\gamma-1}}
    \end{equation}
    \begin{equation}
    \frac{\rho_c}{\rho_1} = \frac{\rho_2}{\rho_1}\left[\frac{1+\frac{\gamma-1}{2}M_2^2}{1+\frac{\gamma-1}{2}M_c^2}\right]^{\frac{1}{\gamma-1}}
    \end{equation}
    \begin{equation}
    C_p = \frac{2}{\gamma M_1^2}\left(\frac{P_c}{P_1}-1\right)
    \end{equation}
    \begin{equation}
    \frac{\Delta s}{R} = \max\left(0,\ -\ln\frac{P_{02}}{P_{01}}\right)
    \end{equation}
    \begin{equation}
    V_{total}' = \sqrt{V_r'(\theta_c)^2+V_\theta'(\theta_c)^2}
    \end{equation}
\end{enumerate}

\subsection*{Step 4: Input Cases and Boundaries}

The three cases of inputs mapped to their boundary limits are:
\begin{itemize}
    \item \textbf{Case 1: Given Shock Angle $\theta_s$}
    \begin{itemize}[nosep]
        \item \textit{Boundary:} $\mu_1 < \theta_s < 90^\circ$, where the lower bound $\mu_1$ is the Mach angle.
        \item \textit{Action:} Solved directly via Step 3.
    \end{itemize}
    \item \textbf{Case 2: Given Cone Angle $\theta_c$}
    \begin{itemize}[nosep]
        \item \textit{Boundary:} $0 < \theta_c \le \theta_{c,max}$, where $\theta_{c,max}$ is the maximum cone angle from Step 1.
        \item \textit{Action:} Guesses $\theta_s$ iteratively using Step 5 to satisfy the target cone angle.
    \end{itemize}
    \item \textbf{Case 3: Given Surface Mach Number $M_c$}
    \begin{itemize}[nosep]
        \item \textit{Boundary:} $M_{c,min} \le M_c < M_1$, where $M_{c,min}$ is the minimum Mach number from Step 1.
        \item \textit{Action:} Guesses $\theta_s$ iteratively using Step 5 to satisfy the target surface Mach number.
    \end{itemize}
\end{itemize}

\subsection*{Step 5: Brent's Method Root-Finding}

Brent's method is a safe, bracketed root-finding algorithm that combines bisection, secant, and inverse quadratic interpolation. In Case 2 and Case 3, the solver drives the following target objective functions to zero:
\begin{itemize}
    \item \textbf{Case 2 objective:} $f(\theta_s) = \theta_c(\theta_s) - \theta_{c,target} = 0$
    \item \textbf{Case 3 objective:} $f(\theta_s) = M_c(\theta_s) - M_{c,target} = 0$
\end{itemize}
over the search interval $[\mu_1 + 10^{-12},\ 90^\circ]$. Brent's method evaluates the target functions (each invoking the central solver of Step 3) until the solution converges to a tolerance of $10^{-15}$.



\section{Program Flow and Architecture}

Conical shock has a unique structure: $\theta_s$ is the central pivot. It is either given directly (Case~1) or found by Brent's root-finding (Cases~2 and~3), then passed to the RK45 Taylor--Maccoll integrator.

\begin{center}
\begin{tikzpicture}[x=1cm,y=1cm,
    case/.style={rectangle, draw=blue!70!black, thick, fill=blue!6, align=center,
                 text width=3.6cm, minimum height=0.85cm, font=\small\sffamily},
    pivot/.style={rectangle, draw=red!70!black, thick, fill=red!7, align=center,
                  text width=1.8cm, minimum height=3.0cm, font=\small\sffamily},
    eng/.style={rectangle, draw=orange!80!black, thick, fill=orange!8, align=center,
                text width=3.0cm, font=\small\sffamily},
    outn/.style={rectangle, draw=green!60!black, thick, fill=green!6, align=center,
                text width=2.8cm, minimum height=0.65cm, font=\small\sffamily},
    arr/.style={-{Stealth[scale=1.0]}, thick, draw=black!70},
    bus/.style={thick, draw=orange!70!black}
]
\node[case]  (c1) at (0, 0.0) {Case 1: $M_1$, $\theta_s$\\(direct)};
\node[case]  (c2) at (0, 1.4) {Case 2: $M_1$, $\theta_c$\\$\to$ Brent's method};
\node[case]  (c3) at (0, 2.8) {Case 3: $M_1$, $M_c$\\$\to$ Brent's method};
\node[pivot] (ts) at (5.0, 1.4) {$\theta_s$\\(pivot)};
\node[eng, minimum height=5.8cm] (eng) at (8.8, 1.4)
    {\textbf{Central Solver}\\[3pt]
    Oblique jump\\[-1pt]
    RK45 TM ODE\\[-1pt]
    Surface detect};
\node[outn] (o0) at (13.2, -0.4) {$\theta_s$, $\theta_c$};
\node[outn] (o1) at (13.2,  0.5) {$M_c$, $M_2$};
\node[outn] (o2) at (13.2,  1.4) {$T_c/T_1$, $T_2/T_1$};
\node[outn] (o3) at (13.2,  2.3) {$P_c/P_1$, $P_2/P_1$};
\node[outn] (o4) at (13.2,  3.2) {$\rho_c/\rho_1$, $\rho_2/\rho_1$};
\node[outn] (o5) at (13.2,  4.1) {$P_{02}/P_{01}$};
\node[outn] (o6) at (13.2,  5.0) {$C_p$, $\Delta s/R$};
\foreach \i in {c1,c2,c3} {\draw[arr] (\i.east) -- (ts.west |- \i);}
\draw[arr] (ts.east) -- (eng.west |- ts);
\coordinate (bx) at ($(eng.east)+(0.5,0)$);
\draw[bus] (eng.east) -- (bx);
\draw[bus] (bx |- o0) -- (bx |- o6);
\foreach \n in {o0,o1,o2,o3,o4,o5,o6} {\draw[arr] (bx |- \n) -- (\n.west);}
\end{tikzpicture}
\end{center}

\noindent Case~1 passes $\theta_s$ directly to the central solver. Cases~2 and~3 use Brent's root-finding over $[\mu_1, 90^\circ]$ to search for $\theta_s$ that matches the target $\theta_c$ or $M_c$, invoking the central solver at each step. The central solver performs an oblique-shock jump at $\theta_s$, then integrates the Taylor--Maccoll ODE with RK45 until $V_\theta'=0$, locating the surface via a nested Brent call (tolerance $10^{-14}$).
